\documentclass[a4paper,10pt]{letter}
\usepackage[margin=0.8in]{geometry}
\usepackage[T1]{fontenc}
\usepackage{hyperref}
\hypersetup{
  hidelinks,
  pdftitle={Revision Cover Letter: Fusion Utility Depends on the Evaluator},
  pdfauthor={Zigao Wu},
  pdfsubject={Revision for Infrared Physics \& Technology}
}
\setlength{\parskip}{0.42em}

\begin{document}
\sloppy
\begin{letter}{The Editor\\Infrared Physics \& Technology}
\opening{Dear Editor,}

We thank you and the reviewers for the constructive assessment of manuscript
INFPHY-D-26-02423, ``Fusion Utility Depends on the Evaluator: A Boundary
Study of Confidence-Shielded Input Selection for Infrared--Visible Object
Detection.'' We also thank you for the invitation to revise by December 25,
2026. We submit a revised manuscript and a separate point-by-point response
for entry in the peer-review platform.

The revision keeps the original visible margin and score depth fixed while
testing whether the full four-candidate library is necessary. On the same four
development evaluator cells, the smallest feasible subset was visible plus
PIAFusion. We then measured the complete pipeline and confirmed the frozen
policy on a later 900-frame LLVIP holdout recorded before scoring. It gained
\(+0.01560\) \([+0.00249,+0.02874]\) AP50 and \(+0.00946\)
\([+0.00113,+0.01769]\) AP50--95 over visible, with positive gains in all
three sequences and six checkpoints. The measured end-to-end median was
\(1.67\) s for the selected pipeline and \(2.52\) s for the full library.

We also retained the boundaries of this result. The same two-candidate policy
was below visible input on VEDAI and on the two independent DroneVehicle runs.
The manuscript therefore presents candidate-library reduction only after
target-specific validation and returns visible input when validation fails.
The later confirmation protocol and its sample amendment were recorded before
candidate or detector scoring, though the policy itself was selected during
revision and the protocol was not externally registered. The measured
pipeline remains too slow for a low-latency claim.

Several reporting and sensitivity gaps raised in review are now closed. We
report the development margin alternatives and the sequence-040 audit values,
the direct no-shield-minus-visible contrast, and the no-detection sentinel
frequency. DroneVehicle now includes confidence-ordered OBB results for both
independent runs and all named comparators. We added a one-at-a-time CDF
sensitivity analysis for qTH, alpha, sigma, and minimum-detail magnitude.
Fixed-CDF AP50 varied across settings, while the selected visible-plus-PIA
policy was unchanged and the full policy changed by at most 0.00018 on test.

The M3FD paired equal-group estimates are now reported explicitly. We also
trained seeds 1 and 2 for both model families with the original 40-epoch recipe
and jointly rescored all six checkpoints. CS-ECF minus visible was positive in
6/6 runs, with a pooled mean of \(+0.00852\) and an equal-group mean of
\(+0.00617\). The magnitude varied across seeds, and the evidence remains
conditional on two fixed frame groups and two model families. We do not present
it as cross-domain confirmation. We also consolidated the limitations and
revised the discussion of prior work to distinguish fixed-detector reporting
from the evaluator- and domain-transfer question studied here.

The clean manuscript, marked revision summary, supplementary material, source
data, and point-by-point responses are provided with this revision. We would
be grateful for re-evaluation.

\closing{Sincerely,}
Zigao Wu\\
North China Electric Power University\\
Baoding, Hebei 071003, China\\
\href{mailto:zgwu@ncepu.edu.cn}{zgwu@ncepu.edu.cn}
\end{letter}
\end{document}
